Like the STM32N6, I followed the \href{https://edit.wpgdadawant.com/uploads/news_file/program/2021/36761/tech_files/program_36761_suggest_other_file.PDF?}{AR0147 datasheet} on how to properly create the image sensor schematic.
Although it wasn't specified in the datasheet, I included common mode chokes on the differential lines like the MIRA220 schematic.
\nl
For the connector to the processor, \href{https://blog.arducam.com/raspberry-pi-camera/connector-type-pinout/}{this guide} was followed on how to properly set the pinout. I set every third connection to ground to ensure a reliable return path for all signals.
\nl
The most important design consideration with this schematic was the requirement of digital and analog ground (DGND/AGND). To understand how to work with different signal planes, the following resources from TI were used: \href{https://www.ti.com/lit/an/slyt499/slyt499.pdf?ts=1756814995128&ref_url=https%3A%2F%2Fwww.google.com%2F}{1} and \href{https://www.ti.com/lit/an/slyt512/slyt512.pdf?ts=1756824651700&ref_url=https%3A%2F%2Fwww.google.com%2F}{2}.
To connect both grounds, I used a simple star connection with a NetTie. I also set DGND as the equivalent ground on the processor PCB (GND=DGND on the connector).
\nl
I selected the \href{https://au.mouser.com/ProductDetail/ECS/ECS-TXO-3225MV-240-TR?qs=d0WKAl%252BL4KYXn7UuwSFJGQ%3D%3D&srsltid=AfmBOoqlFJCYyqd6MzLT63_Pw_wqXWR5Bv4SHlYdyhoVLS8l8SGV9TkK}{ECS-TXO-3225MV-240-TR} as the external PLL oscillator for the image sensor. This component can be operated on 3V3, however, I chose 2V8 so the oscillator will only start when the sensor's specific power domains are provided and doesn't 'burn' power when the sensor is not in use.
Page 22 of the sensor datasheet specifies that the tEXTCLK signal must have a rise/fall of between 0.2ns and 12.5ns (12.5ns extrapolated from 0.3*tEXTCLK). These requirements are met by the selected oscillator since the oscillator datasheet says it has a rise/fall time of 5$\pm$0.5ns. Since the generated signal is a clipped sine wave, I can estimate the duty cycle is 50\% since it's symmetric.
